\section{Conclusion}
\label{sec:Con}

This study evaluated a minimal runtime verification pipeline for automotive ECU telemetry replay, using LARVA-driven monitoring over four properties (\texttt{voltageOutOfRange}, \texttt{coolantTooHigh}, \texttt{throttleSpike}, \texttt{loadMapMismatch}). The objective was to quantify monitoring overhead and detection behaviour under controlled baseline and abnormal traces.

Findings indicate that: (i) injected abnormal classes are detected consistently when RV is enabled and not reported when RV is disabled; (ii) on the stock generated monitor, nominal overhead is moderate (6.43\%) but abnormal-path overhead is extreme (3927.05\%) before diagnostic optimisations; and (iii) threshold tuning reduces false positives without collapsing abnormal detection. The controlled three-variant comparison showed that abnormal overhead is mainly driven by generated bad-state diagnostics: 3927.05\% (original generated) decreased to 683.79\% (logging disabled) and then to 9.39\% after removing stacktrace-based bad-state string construction, with CPU-time following the same trend. The hypothesis is \textbf{partially supported}: useful coverage is achievable with bounded nominal overhead, and worst-case overhead can be reduced through targeted monitor-code optimisation.

Methodological limitations should be acknowledged. Evaluation is trace-driven and server-side, not hardware-in-the-loop on target automotive microcontrollers; measured overhead therefore reflects server-path cost rather than on-device monitor cost. The selected property set is intentionally small, and the abnormal dataset is synthetic (though controlled), so broader production diversity is not fully represented.

Future work should move to ECU-representative hardware, expand validation with recorded CAN/FlexRay and hardware-in-the-loop campaigns, and integrate trust primitives (attestation and hardware trust anchors) with monitor lifecycle management~\cite{ekert2021cybersecurity,plappert2023hta,lorych2024hta}.